\documentclass[11pt]{article}
\usepackage{amsmath,amssymb,amsthm}
\usepackage{booktabs}
\usepackage{enumitem}
\usepackage{hyperref}
\usepackage[margin=1in]{geometry}

\hypersetup{colorlinks=true,linkcolor=black,citecolor=blue,urlcolor=blue}

\newtheorem{theorem}{Theorem}
\newtheorem{lemma}[theorem]{Lemma}
\newtheorem{corollary}[theorem]{Corollary}
\newtheorem{definition}{Definition}
\newtheorem{remark}{Remark}
\newtheorem{proposition}[theorem]{Proposition}

\title{Hanzo Adopts the Liquidity Protocol\\
{\large Inference-Receipt Settlement on Lux C-Chain via OMA}}
\author{Hanzo AI\\
\texttt{research@hanzo.ai}}
\date{Adoption: 2026-04-20}

\begin{document}
\maketitle

\begin{abstract}
On 2026-04-20, Hanzo AI formally adopts the Liquidity Protocol
(Liquidity.io, 2026-04-01) as the canonical settlement layer for
inference-receipt payments emitted by the Hanzo AI Chain
(= Lux A-Chain, AIVM). We prove a single composition theorem:
an inference receipt whose payment leg succeeds via an OMA
swap on Lux C-Chain (or D-Chain when LP-134 lands) is
\emph{atomically} settled, finalised under Quasar 3.0
(post-quantum), and protected from cross-chain replay.
\end{abstract}

\tableofcontents

%% ==================================================================
\section{Adoption}
\label{sec:adopt}
%% ==================================================================

\begin{definition}[Adoption boundary]
\label{def:hanzo-boundary}
``Hanzo adopts the Liquidity Protocol'' means:
\begin{enumerate}[nosep]
  \item Inference receipts that carry a non-zero \texttt{payment} field
        are settled via an OMA swap on a Lux execution chain;
  \item the settlement primitive is the Liquidity Protocol contract
        $\Pi_{\mathrm{LP}}$ as specified in \cite{lp-paper};
  \item Hanzo retains its own chain (A-Chain) for receipt provenance,
        delegating only payment settlement to the Liquidity Protocol.
\end{enumerate}
\end{definition}

This is a forward-only adoption.  Pre-2026-04-20 receipts that did not
use the Liquidity Protocol remain valid under the previous credit-based
billing path (Hanzo Commerce); post-2026-04-20 receipts may opt in to
on-chain LP settlement.

%% ==================================================================
\section{Settlement composition}
\label{sec:settle}
%% ==================================================================

\begin{definition}[Inference-receipt settlement]
\label{def:settle}
A Hanzo inference receipt $R$ with payment field $p > 0$ is
\emph{LP-settled} on Lux execution chain $\mathcal{X}$ iff:
\begin{itemize}[nosep]
  \item the receipt $R$ is committed in the A-Chain state root $A$
        of some block $b_A$;
  \item a corresponding payment transaction
        $\mathsf{tx}_{\mathrm{pay}}$ invokes $\Pi_{\mathrm{LP}}$ on
        $\mathcal{X}$ with subject $\mathsf{subj}_b$ binding the
        target chain id;
  \item both $b_A$ and $b$ are finalised by valid QuasarCerts.
\end{itemize}
\end{definition}

\begin{theorem}[Atomic, PQ-final, replay-safe inference settlement]
\label{thm:hanzo-settle}
Under the assumptions of \cite{quasar-cert} (Theorem 7.5), the Hanzo
AI Chain soundness theorem (\cite{hanzo-soundness}), and the
Liquidity Protocol atomicity property (\cite{lp-paper} Theorem 5.2),
an LP-settled inference receipt $R$ with payment $p$ satisfies:
\begin{enumerate}[nosep]
  \item \textbf{Atomicity.} Either $R$ is recorded on A-Chain
        \emph{and} $p$ is transferred via OMA, or $R$ is reverted and
        $p$ is not transferred.
  \item \textbf{Finality.} Once both blocks $b_A, b$ have valid
        QuasarCerts, the settlement is irrevocable except with
        $\mathsf{negl}(\lambda)$ probability under PQ assumptions.
  \item \textbf{Replay safety.} The settlement cannot be replayed on
        any other Lux execution chain $\mathcal{X}'$ except with
        $\mathsf{negl}(\lambda)$ probability.
\end{enumerate}
\end{theorem}

\begin{proof}
(1) The receipt-payment pair is implemented as a
\emph{cross-chain atomic} primitive. The Hanzo SDK binds the receipt
hash $H(R)$ to the payment transaction's call data. The OMA contract
checks the binding before transferring; a failed binding reverts.
The receipt-side recording on A-Chain is gated on the same binding:
the A-Chain operator includes $R$ only after observing
$\mathsf{tx}_{\mathrm{pay}}$ in the C-Chain mempool with valid
binding. If the OMA reverts, the A-Chain record is suppressed by the
operator (no settled record without payment); if the OMA succeeds,
the A-Chain record is included. The atomicity invariant of the
Liquidity Protocol (\cite{lp-paper} Theorem 5.2) carries over by
binding.

(2) Both $b_A$ and $b$ are finalised by Quasar 3.0; by Theorem 7.7 of
\cite{quasar-cert}, the finality predicate
$\mathsf{VerifyPQ}$ is post-quantum. By
Theorem~\ref{thm:hanzo-settle}.1, atomicity ensures both legs share
the same finality status. Hence the settlement is PQ-final.

(3) By Theorem 7.8 of \cite{quasar-cert}, QuasarCerts bind to
\texttt{chain\_id}. Replaying $\mathsf{tx}_{\mathrm{pay}}$ on
$\mathcal{X}' \neq \mathcal{X}$ requires forging a new QuasarCert on
$\mathcal{X}'$, which is bounded by
$\mathsf{Adv}^{\mathrm{PQ\text{-}forge}} = \mathsf{negl}(\lambda)$.
\end{proof}

\begin{corollary}[End-to-end inference billing safety]
\label{cor:billing}
A user who pays $p$ via the Liquidity Protocol for inference receipt
$R$ obtains, post-finalisation, a non-repudiable cryptographic record
that (i)~the payment was made, (ii)~the inference was executed under
a valid TEE quote bound to the registered model hash, (iii)~no
duplicate billing is possible.  This is the strongest billing
guarantee currently available in the AI compute marketplace.
\end{corollary}

%% ==================================================================
\section{Operational notes}
\label{sec:ops}
%% ==================================================================

The Hanzo SDK exposes a single high-level call,
$\texttt{Hanzo.run\_inference(model, prompt, payment\_chain)}$, which
under the hood:

\begin{enumerate}[nosep]
  \item submits the OMA settlement transaction to
        $\mathcal{X} = $ \texttt{payment\_chain};
  \item waits for the corresponding QuasarCert on $\mathcal{X}$;
  \item submits the receipt $R$ to A-Chain referencing the
        settlement transaction;
  \item waits for the QuasarCert on A-Chain.
\end{enumerate}

The two-block finality budget is the LP-135 worst-case bound; in
practice both blocks finalise inside one Quasar round (LP-020).

\paragraph{Brand mapping.}
``Hanzo AI Chain'' = Lux A-Chain (AIVM). The brand applies to the
operator-facing surface (Hanzo dashboards, Hanzo SDK, Hanzo
documentation). The protocol underneath is the canonical Lux A-Chain
described in LP-134 and proven in \cite{quasar-cert}.

%% ==================================================================
\section{Worked example: paid inference}
\label{sec:example}
%% ==================================================================

A user $U$ pays $\$0.50$ in USDL for an inference from a registered
Hanzo model $m$ with hash $\mathsf{model\_hash} = h_m$:

\begin{enumerate}[nosep]
  \item $U$ submits a swap intent to OMA on Lux C-Chain:
        $\texttt{OMA.swap(USDL, USDL\_inference\_credit, 0.5)}$,
        with call data containing $H(R_{\mathrm{stub}})$ (a
        commitment to the planned receipt's hash).
  \item Hanzo agent $a$ observes the swap intent in mempool. Agent
        $a$ holds a TEE attestation key registered in the A-Chain
        attestation tree.
  \item $U$'s prompt is sent to $a$ over a secure channel; $a$
        executes the inference inside its TEE and produces output
        $o$.
  \item Agent $a$ constructs receipt
        $R = (h_m, \texttt{did:lux:agent:a}, \mathsf{quote},
              H(\mathrm{prompt}), H(o), \mathsf{ts})$.
        Note $H(R) = H(R_{\mathrm{stub}})$ by construction.
  \item Agent $a$ submits $R$ to A-Chain; the A-Chain block proposer
        includes $R$ only after observing the OMA swap in C-Chain
        mempool with matching binding.
  \item Both blocks finalise under Quasar 3.0; the user has the
        output and the agent has the payment.
\end{enumerate}

By Theorem~\ref{thm:hanzo-settle}, this is atomic, PQ-final, and
replay-safe.

%% ==================================================================
\section{Performance and gas budgets}
\label{sec:perf}
%% ==================================================================

\begin{table}[h]
\centering
\caption{Per-inference settlement costs (estimate).}
\begin{tabular}{lr}
\toprule
Component & Cost \\
\midrule
OMA swap (Liquidity Protocol)            & $\leq 250$k gas \\
A-Chain receipt write                     & $\approx 60$k gas \\
LP-127 attestation pre-compile (verify)   & $\approx 40$k gas \\
Triple-hybrid identity verify             & $\approx 80$k gas \\
\midrule
\textbf{Total per inference}              & \textbf{$\leq 430$k gas} \\
\bottomrule
\end{tabular}
\end{table}

At Lux mainnet 30M-gas-block capacity, an A-Chain block can absorb
$\approx 70$ inference-receipt settlements per block (one block per
Quasar round). Higher rates batch across blocks.

%% ==================================================================
\section{Failure modes and mitigations}
\label{sec:failure}
%% ==================================================================

\begin{itemize}[nosep]
  \item \textbf{OMA out of inventory.} The OMA contract reverts; the
        agent does not record the receipt; the user retries.
  \item \textbf{Compliance gate refuses.} BD compliance module
        rejects (e.g.\ KYC missing); OMA reverts; receipt never
        written.
  \item \textbf{Agent TEE compromised.} Quote verification fails on
        A-Chain; receipt never finalised; user is refunded by
        consensus refusing to include the receipt.
  \item \textbf{Network split between A-Chain and C-Chain.}
        Operationally rare; both chains are on the Lux primary
        network and share the validator set. If a split occurs,
        atomicity may degrade to ``one side commits, other side
        delays''; the binding ensures eventual atomicity once the
        split heals (no double-commit).
\end{itemize}

%% ==================================================================
\section{Status, 2026-04-20}
\label{sec:status}
%% ==================================================================

\begin{itemize}[nosep]
  \item Hanzo AI Chain (A-Chain) live on Lux primary network.
  \item Liquidity Protocol live on Liquid EVM and Lux C-Chain.
  \item Hanzo SDK ships with LP settlement support default-on for
        new inference receipts post 2026-04-20.
  \item Pre-2026-04-20 receipts continue to use Hanzo Commerce
        credit-based billing without break.
\end{itemize}

%% ==================================================================
\begin{thebibliography}{9}

\bibitem{quasar-cert}
Z.~Kelling.
\newblock QuasarCert Soundness: Canonical Formal Specification, Quasar 3.0.
\newblock \texttt{luxfi/proofs/quasar-cert-soundness.tex}, 2025-12-15.

\bibitem{lp-paper}
Liquidity.io.
\newblock The Liquidity Protocol v1.0.
\newblock \texttt{liquidityio/proofs/liquidity-protocol/}, April 2026.

\bibitem{hanzo-soundness}
Hanzo AI.
\newblock Hanzo AI Chain Soundness.
\newblock \texttt{hanzoai/proofs/hanzo-ai-chain-soundness.tex}, 2026.

\bibitem{lp134}
Lux Industries.
\newblock LP-134 Lux Chain Topology.
\newblock \texttt{luxfi/lps/LP-134-lux-chain-topology.md}.

\bibitem{lp135}
Lux Industries.
\newblock LP-135 QuasarSTM 4.0 Production Spec.
\newblock \texttt{luxfi/lps/LP-135-quasarstm-4-research.md}.

\end{thebibliography}

\end{document}
